\hypertarget{appendix-215c2570199b}{}\label{referencestart:appendix-215c2570199b}
\section{Light spectra and radiometry}\label{appendix-215c2570199b-light-spectra-and-radiometry}

\textbf{Curriculum depth:} Developed · \textbf{First seen:} XR01 · \textbf{Latest development:} University courses

\subsection{Definition and mechanism}\label{appendix-215c2570199b-definition-and-mechanism}

Light models distinguish wave propagation, photon energy and geometrical rays. Radiometry describes physical flux/radiance/irradiance; photometry adds standardized visual weighting. RGB requires an interpretation and cannot identify a unique spectrum.

\subsection{Prerequisites}\label{appendix-215c2570199b-prerequisites}

No formal graph prerequisite. Begin with the stated units and model; this is an independent-series foundation.

\subsection{Used by}\label{appendix-215c2570199b-used-by}

\begin{itemize}
\tightlist
\item
  \hyperref[appendix-96da0b153b77]{Reflectance and rendering equation}
\item
  \hyperref[appendix-7281c15ac44a]{Retinal encoding and receptive fields}
\item
  \hyperref[appendix-f90cee1be7a4]{Optical focus and binocular conflict}
\item
  \hyperref[appendix-a9f83ad40792]{Color encoding and linear-light operations}
\item
  \hyperref[appendix-4815cd87ab8b]{Probability density and change of measure}
\item
  \hyperref[appendix-89b2d5b7c9e1]{Surface orientation from image brightness}
\end{itemize}

\subsection{Recall and next depth}\label{appendix-215c2570199b-recall-and-next-depth}

Explain W, W/m², W/(m²·sr), lux and nits. Pending: spectral/colorimetric calibration. Written coverage is not a claim of learner mastery or accepted numbered-curriculum progress.

\subsection{Source and navigation}\label{appendix-215c2570199b-source-and-navigation}

\href{../07-Vision-XR-and-Playable-Worlds/main.pdf\#nameddest=04-visual-computing-ii-ch042-light-must-become-a-quantity-before-it-becomes-a-color}{Expanded teaching explanation (Vision, XR, and Playable Worlds, Vision, XR, and Playable Worlds, Chapter 9)} · \hyperref[appendix-fd090d982a15-xr01---independent-seeing-and-representation-source]{Independent coverage} · \hyperref[appendix-4ad47bdd269a]{Dependency Map}

\subsection{University course development}\label{appendix-215c2570199b-university-course-development}

Radiometry distinguishes total radiant flux \(\Phi\) in watts, irradiance \(E\) in watts per square metre, radiant intensity \(I\) in watts per steradian, and radiance \(L\) in watts per square metre per steradian measured with projected source area. A spectrum distributes the energy over wavelength; a three-channel RGB value is a measurement in a chosen color system and does not uniquely identify that spectrum. Solid angle measures directional extent rather than physical area. Consequently an area-light sample requires the projected-area and inverse-square Jacobian when expressed as a directional sample. Free-space radiance remains constant along an unobstructed ray, while irradiance from a point-like source changes with geometry. Photometry instead applies standardized visual weighting: luminous flux uses lumens, illuminance lux, and luminous intensity candela. These distinctions prevent treating brightness, exposure, flux and radiance as interchangeable scalar labels.

\href{../03-Geometry-Cameras-and-Image-Formation/main.pdf\#nameddest=03-visual-computing-i-ch004-light-measurement-and-radiometry}{Light measurement and radiometry (Geometry, Cameras, and Image Formation, Geometry, Cameras, and Image Formation, Chapter 11)}
